% The lab's paper: every task side by side. The prose follows docs/writing.md.
% Every number, table and figure is written by `lab paper` into generated/,
% from every model's committed results, and read back with \V{key} and
% \block{name}. A task's values and blocks stand under its name.
\documentclass[11pt]{article}
% The preamble every paper of the lab shares: \input{../../../papers/preamble}
% from models/<model>/paper/, % The preamble every paper of the lab shares: \input{../../../papers/preamble}
% from models/<model>/paper/, % The preamble every paper of the lab shares: \input{../../../papers/preamble}
% from models/<model>/paper/, \input{../preamble} from papers/<name>/.
% Boards come from chessboard and xskak, plots from pgfplots. A paper prints no
% number by hand: `lab paper` writes every number, table and figure under the
% paper's generated/ folder, and paper.tex says what they mean.
\usepackage[LSBC4,T1]{fontenc}
\usepackage[utf8]{inputenc}
\usepackage{lmodern}
\usepackage[margin=1in]{geometry}
\usepackage{microtype}
\usepackage{amsmath}
\usepackage{booktabs}
\usepackage{tabularx}
\usepackage{array}
\usepackage{xcolor}
\usepackage[obeyspaces,spaces]{url}
\usepackage{xskak}
\usepackage{pgfplots}
\usepgfplotslibrary{groupplots}
\usetikzlibrary{positioning,arrows.meta,calc,fit}
\pgfplotsset{compat=1.18}
\usepackage{caption}
\usepackage{subcaption}
\usepackage{enumitem}
\usepackage[section]{placeins}
\usepackage[hidelinks]{hyperref}

% A generated value by its key: generated/values.tex sets them, \V{key} prints
% one, and a key nobody set stops the build.
\makeatletter
\newcommand{\setvalue}[2]{\expandafter\gdef\csname labvalue@#1\endcsname{#2}}
\newcommand{\V}[1]{\@ifundefined{labvalue@#1}%
  {\PackageError{lab}{No generated value named '#1'}{Run `lab paper` for this model again.}}%
  {\csname labvalue@#1\endcsname}}
\makeatother

% Chart colours. A run keeps its colour in every figure of every paper.
\definecolor{labBlue}{HTML}{2A78D6}
\definecolor{labOrange}{HTML}{EB6834}
\definecolor{labAqua}{HTML}{1BAF7A}
\definecolor{labYellow}{HTML}{EDA100}
\definecolor{labMagenta}{HTML}{E87BA4}
\definecolor{labGreen}{HTML}{008300}
\definecolor{labViolet}{HTML}{6250D6}
\definecolor{labGray}{HTML}{898781}
\definecolor{labInk}{HTML}{0B0B0B}
\definecolor{labSecondary}{HTML}{52514E}
\definecolor{labGrid}{HTML}{E1E0D9}
\definecolor{labAxis}{HTML}{C3C2B7}
\definecolor{labWash}{HTML}{CDE2FB}

\pgfplotsset{
  lab/.style={
    font=\sffamily\footnotesize,
    axis line style={labAxis},
    tick style={labAxis},
    tick label style={labSecondary},
    label style={labSecondary},
    title style={labInk, font=\sffamily\small\bfseries, align=left},
    grid style={labGrid, line width=0.3pt},
    major grid style={labGrid, line width=0.3pt},
    axis lines*=left,
    tick align=outside,
    clip=false,
    legend style={draw=none, fill=none, font=\sffamily\footnotesize, text=labSecondary,
      cells={anchor=west}, legend columns=-1, column sep=6pt, at={(0,1.03)}, anchor=south west},
    every axis plot/.append style={line width=1.1pt, line join=round, line cap=round},
    scaled ticks=false,
    /pgf/number format/1000 sep={,},
  },
  lab marks/.style={mark=*, mark size=1.9pt, mark options={draw=white, line width=0.6pt}},
}

\setchessboard{boardfontencoding=LSBC4, boardfontfamily=skaknew, blackfieldcolor=black!18,
  whitefieldcolor=white, setfontcolors, showmover=false, label=true, labelfontsize=6pt,
  boardfontsize=12pt}

% A move or line in UCI, a position in FEN (breakable at a slash), and a run's name.
\newcommand{\uci}[1]{\texttt{#1}}
\newcommand{\fen}[1]{{\footnotesize\path{#1}}}
\newcommand{\run}[1]{\texttt{#1}}
\newcolumntype{L}{>{\raggedright\arraybackslash}X}
\newcolumntype{R}{>{\raggedleft\arraybackslash}X}

\setlength{\parskip}{0.4em}
\setlength{\parindent}{0pt}
\captionsetup{font=small, labelfont=bf}
 from papers/<name>/.
% Boards come from chessboard and xskak, plots from pgfplots. A paper prints no
% number by hand: `lab paper` writes every number, table and figure under the
% paper's generated/ folder, and paper.tex says what they mean.
\usepackage[LSBC4,T1]{fontenc}
\usepackage[utf8]{inputenc}
\usepackage{lmodern}
\usepackage[margin=1in]{geometry}
\usepackage{microtype}
\usepackage{amsmath}
\usepackage{booktabs}
\usepackage{tabularx}
\usepackage{array}
\usepackage{xcolor}
\usepackage[obeyspaces,spaces]{url}
\usepackage{xskak}
\usepackage{pgfplots}
\usepgfplotslibrary{groupplots}
\usetikzlibrary{positioning,arrows.meta,calc,fit}
\pgfplotsset{compat=1.18}
\usepackage{caption}
\usepackage{subcaption}
\usepackage{enumitem}
\usepackage[section]{placeins}
\usepackage[hidelinks]{hyperref}

% A generated value by its key: generated/values.tex sets them, \V{key} prints
% one, and a key nobody set stops the build.
\makeatletter
\newcommand{\setvalue}[2]{\expandafter\gdef\csname labvalue@#1\endcsname{#2}}
\newcommand{\V}[1]{\@ifundefined{labvalue@#1}%
  {\PackageError{lab}{No generated value named '#1'}{Run `lab paper` for this model again.}}%
  {\csname labvalue@#1\endcsname}}
\makeatother

% Chart colours. A run keeps its colour in every figure of every paper.
\definecolor{labBlue}{HTML}{2A78D6}
\definecolor{labOrange}{HTML}{EB6834}
\definecolor{labAqua}{HTML}{1BAF7A}
\definecolor{labYellow}{HTML}{EDA100}
\definecolor{labMagenta}{HTML}{E87BA4}
\definecolor{labGreen}{HTML}{008300}
\definecolor{labViolet}{HTML}{6250D6}
\definecolor{labGray}{HTML}{898781}
\definecolor{labInk}{HTML}{0B0B0B}
\definecolor{labSecondary}{HTML}{52514E}
\definecolor{labGrid}{HTML}{E1E0D9}
\definecolor{labAxis}{HTML}{C3C2B7}
\definecolor{labWash}{HTML}{CDE2FB}

\pgfplotsset{
  lab/.style={
    font=\sffamily\footnotesize,
    axis line style={labAxis},
    tick style={labAxis},
    tick label style={labSecondary},
    label style={labSecondary},
    title style={labInk, font=\sffamily\small\bfseries, align=left},
    grid style={labGrid, line width=0.3pt},
    major grid style={labGrid, line width=0.3pt},
    axis lines*=left,
    tick align=outside,
    clip=false,
    legend style={draw=none, fill=none, font=\sffamily\footnotesize, text=labSecondary,
      cells={anchor=west}, legend columns=-1, column sep=6pt, at={(0,1.03)}, anchor=south west},
    every axis plot/.append style={line width=1.1pt, line join=round, line cap=round},
    scaled ticks=false,
    /pgf/number format/1000 sep={,},
  },
  lab marks/.style={mark=*, mark size=1.9pt, mark options={draw=white, line width=0.6pt}},
}

\setchessboard{boardfontencoding=LSBC4, boardfontfamily=skaknew, blackfieldcolor=black!18,
  whitefieldcolor=white, setfontcolors, showmover=false, label=true, labelfontsize=6pt,
  boardfontsize=12pt}

% A move or line in UCI, a position in FEN (breakable at a slash), and a run's name.
\newcommand{\uci}[1]{\texttt{#1}}
\newcommand{\fen}[1]{{\footnotesize\path{#1}}}
\newcommand{\run}[1]{\texttt{#1}}
\newcolumntype{L}{>{\raggedright\arraybackslash}X}
\newcolumntype{R}{>{\raggedleft\arraybackslash}X}

\setlength{\parskip}{0.4em}
\setlength{\parindent}{0pt}
\captionsetup{font=small, labelfont=bf}
 from papers/<name>/.
% Boards come from chessboard and xskak, plots from pgfplots. A paper prints no
% number by hand: `lab paper` writes every number, table and figure under the
% paper's generated/ folder, and paper.tex says what they mean.
\usepackage[LSBC4,T1]{fontenc}
\usepackage[utf8]{inputenc}
\usepackage{lmodern}
\usepackage[margin=1in]{geometry}
\usepackage{microtype}
\usepackage{amsmath}
\usepackage{booktabs}
\usepackage{tabularx}
\usepackage{array}
\usepackage{xcolor}
\usepackage[obeyspaces,spaces]{url}
\usepackage{xskak}
\usepackage{pgfplots}
\usepgfplotslibrary{groupplots}
\usetikzlibrary{positioning,arrows.meta,calc,fit}
\pgfplotsset{compat=1.18}
\usepackage{caption}
\usepackage{subcaption}
\usepackage{enumitem}
\usepackage[section]{placeins}
\usepackage[hidelinks]{hyperref}

% A generated value by its key: generated/values.tex sets them, \V{key} prints
% one, and a key nobody set stops the build.
\makeatletter
\newcommand{\setvalue}[2]{\expandafter\gdef\csname labvalue@#1\endcsname{#2}}
\newcommand{\V}[1]{\@ifundefined{labvalue@#1}%
  {\PackageError{lab}{No generated value named '#1'}{Run `lab paper` for this model again.}}%
  {\csname labvalue@#1\endcsname}}
\makeatother

% Chart colours. A run keeps its colour in every figure of every paper.
\definecolor{labBlue}{HTML}{2A78D6}
\definecolor{labOrange}{HTML}{EB6834}
\definecolor{labAqua}{HTML}{1BAF7A}
\definecolor{labYellow}{HTML}{EDA100}
\definecolor{labMagenta}{HTML}{E87BA4}
\definecolor{labGreen}{HTML}{008300}
\definecolor{labViolet}{HTML}{6250D6}
\definecolor{labGray}{HTML}{898781}
\definecolor{labInk}{HTML}{0B0B0B}
\definecolor{labSecondary}{HTML}{52514E}
\definecolor{labGrid}{HTML}{E1E0D9}
\definecolor{labAxis}{HTML}{C3C2B7}
\definecolor{labWash}{HTML}{CDE2FB}

\pgfplotsset{
  lab/.style={
    font=\sffamily\footnotesize,
    axis line style={labAxis},
    tick style={labAxis},
    tick label style={labSecondary},
    label style={labSecondary},
    title style={labInk, font=\sffamily\small\bfseries, align=left},
    grid style={labGrid, line width=0.3pt},
    major grid style={labGrid, line width=0.3pt},
    axis lines*=left,
    tick align=outside,
    clip=false,
    legend style={draw=none, fill=none, font=\sffamily\footnotesize, text=labSecondary,
      cells={anchor=west}, legend columns=-1, column sep=6pt, at={(0,1.03)}, anchor=south west},
    every axis plot/.append style={line width=1.1pt, line join=round, line cap=round},
    scaled ticks=false,
    /pgf/number format/1000 sep={,},
  },
  lab marks/.style={mark=*, mark size=1.9pt, mark options={draw=white, line width=0.6pt}},
}

\setchessboard{boardfontencoding=LSBC4, boardfontfamily=skaknew, blackfieldcolor=black!18,
  whitefieldcolor=white, setfontcolors, showmover=false, label=true, labelfontsize=6pt,
  boardfontsize=12pt}

% A move or line in UCI, a position in FEN (breakable at a slash), and a run's name.
\newcommand{\uci}[1]{\texttt{#1}}
\newcommand{\fen}[1]{{\footnotesize\path{#1}}}
\newcommand{\run}[1]{\texttt{#1}}
\newcolumntype{L}{>{\raggedright\arraybackslash}X}
\newcolumntype{R}{>{\raggedleft\arraybackslash}X}

\setlength{\parskip}{0.4em}
\setlength{\parindent}{0pt}
\captionsetup{font=small, labelfont=bf}

% Written by `lab paper`; do not edit by hand.
\waitingtrue
\changehandtrue
\changeepochstrue
\changestartfalse
\gamestrue
\featurestrue
\studiestrue
\setvalue{blunder-hazard.best.calibration.groups}{ten}
\setvalue{blunder-hazard.best.calibration.widest-gap}{0.8}
\setvalue{blunder-hazard.best.score}{0.2666}
\setvalue{blunder-hazard.best.top.1.lift}{7.1}
\setvalue{blunder-hazard.best.top.1.rate}{48.4}
\setvalue{blunder-hazard.context.lift.share}{65}
\setvalue{blunder-hazard.data.games.test}{2{,}944}
\setvalue{blunder-hazard.data.month}{October 2026}
\setvalue{blunder-hazard.data.rows}{698{,}013}
\setvalue{blunder-hazard.floor.constant}{6.7\%}
\setvalue{blunder-hazard.floor.score}{0.0685}
\setvalue{blunder-hazard.recipe.games}{30{,}000}
\setvalue{blunder-hazard.recipe.speeds}{rapid}
\setvalue{blunder-hazard.reference.score}{0.2533}
\setvalue{lab.board.ahead.names}{puzzle difficulty, puzzle difficulty unsolved and think time}
\setvalue{lab.board.ahead.of}{three of the six models}
\setvalue{lab.board.behind.names}{blunder hazard and tactic found}
\setvalue{lab.board.behind.of}{two of the six models}
\setvalue{lab.board.level.names}{practical chances}
\setvalue{lab.board.level.of}{one of the six models}
\setvalue{lab.change.epochs.net.better.of}{five of the eight pairs}
\setvalue{lab.change.epochs.net.level}{three}
\setvalue{lab.change.epochs.over.transformer-medium-hand.gain.most}{+4.0}
\setvalue{lab.change.epochs.over.transformer-medium-hand.said}{better beyond the noise in two of the four models (puzzle difficulty and think time) and inside the noise in two of the four models (blunder hazard and tactic found)}
\setvalue{lab.change.epochs.over.transformer-medium.count}{four}
\setvalue{lab.change.epochs.over.transformer-medium.gain.most}{+9.4}
\setvalue{lab.change.epochs.over.transformer-medium.said}{better beyond the noise in three of the four models (puzzle difficulty, tactic found and think time) and inside the noise in one of the four models (blunder hazard)}
\setvalue{lab.change.hand.over.gbm-context.said}{better beyond the noise in all five models (blunder hazard, practical chances, puzzle difficulty, tactic found and think time)}
\setvalue{lab.change.hand.over.transformer-medium.models}{blunder hazard, practical chances, puzzle difficulty, puzzle difficulty unsolved, tactic found and think time}
\setvalue{lab.change.hand.over.transformer-medium.said}{better beyond the noise in five of the six models (blunder hazard, puzzle difficulty, puzzle difficulty unsolved, tactic found and think time) and inside the noise in one of the six models (practical chances)}
\setvalue{lab.context.share.least}{35}
\setvalue{lab.context.share.least.model}{tactic found}
\setvalue{lab.context.share.most}{97}
\setvalue{lab.context.share.most.model}{practical chances}
\setvalue{lab.cost.seconds.networks}{4.6 h}
\setvalue{lab.cost.seconds.others}{4 min}
\setvalue{lab.cost.seconds.studies}{2.1 h}
\setvalue{lab.features.in.blunder-hazard.most.gain}{$-$18.7}
\setvalue{lab.features.in.blunder-hazard.most.group}{material}
\setvalue{lab.features.in.practical-chances.most.gain}{$-$2.5}
\setvalue{lab.features.in.practical-chances.most.group}{material}
\setvalue{lab.features.in.puzzle-difficulty.most.gain}{$-$21.0}
\setvalue{lab.features.in.puzzle-difficulty.most.group}{the line}
\setvalue{lab.features.in.tactic-found.most.gain}{$-$16.8}
\setvalue{lab.features.in.tactic-found.most.group}{the first move}
\setvalue{lab.features.in.think-time.most.gain}{$-$9.9}
\setvalue{lab.features.in.think-time.most.group}{material}
\setvalue{lab.games.gbm-hand.half.gain.least}{$-$3.8}
\setvalue{lab.games.gbm-hand.half.gain.least.model}{tactic found}
\setvalue{lab.games.gbm-hand.half.said}{inside the noise in two of the five models (blunder hazard and practical chances) and worse beyond the noise in three of the five models (puzzle difficulty, tactic found and think time)}
\setvalue{lab.games.gbm-hand.quarter.said}{inside the noise in one of the five models (practical chances) and worse beyond the noise in four of the five models (blunder hazard, puzzle difficulty, tactic found and think time)}
\setvalue{lab.games.transformer-medium-hand.half.gain.least}{$-$9.0}
\setvalue{lab.games.transformer-medium-hand.half.gain.least.model}{tactic found}
\setvalue{lab.games.transformer-medium-hand.half.said}{worse beyond the noise in all five models (blunder hazard, practical chances, puzzle difficulty, tactic found and think time)}
\setvalue{lab.games.transformer-medium-hand.quarter.said}{worse beyond the noise in all five models (blunder hazard, practical chances, puzzle difficulty, tactic found and think time)}
\setvalue{lab.hidden.names}{puzzle difficulty unsolved}
\setvalue{lab.ladder.at-last-epoch.of}{twelve of the 24 networks}
\setvalue{lab.models.count}{six}
\setvalue{lab.network.ahead.names}{blunder hazard, puzzle difficulty, puzzle difficulty unsolved, tactic found and think time}
\setvalue{lab.network.ahead.of}{five of the six models}
\setvalue{lab.network.ahead.runs}{\run{transformer-medium-hand}, \run{transformer-medium-hand-long} and \run{conv-small}}
\setvalue{lab.network.behind.of}{none of the six models}
\setvalue{lab.network.gain.whole.most}{16}
\setvalue{lab.network.gain.whole.most.model}{puzzle difficulty}
\setvalue{lab.network.level.names}{practical chances}
\setvalue{lab.network.level.of}{one of the six models}
\setvalue{lab.networks.count}{32}
\setvalue{lab.results.date}{22 November 2026}
\setvalue{lab.runs.count}{52}
\setvalue{lab.small.ahead.of}{two of the six models}
\setvalue{lab.small.behind.of}{two of the six models}
\setvalue{lab.small.level.of}{two of the six models}
\setvalue{lab.split.test.percent}{10}
\setvalue{lab.split.valid.percent}{5}
\setvalue{lab.studies.count}{56}
\setvalue{lab.waiting.names}{position evaluation}
\setvalue{lab.whole.count}{five}
\setvalue{method.blunder-bar}{20}
\setvalue{method.resamples}{1{,}000}
\setvalue{practical-chances.context.lift.share}{97}
\setvalue{practical-chances.grid.rating-win.-1100.80+.happened}{0.760}
\setvalue{practical-chances.grid.rating-win.2000+.80+.happened}{0.901}
\setvalue{practical-chances.overall.engine-curve.said}{0.525}
\setvalue{practical-chances.overall.happened}{0.502}
\setvalue{practical-chances.pair.gbm-context.engine-curve.gain}{+0.01136}
\setvalue{practical-chances.pair.gbm-context.engine-curve.gain.high}{+0.01369}
\setvalue{practical-chances.pair.gbm-context.engine-curve.gain.low}{+0.00892}
\setvalue{practical-chances.run.engine-curve.lift.share}{80}
\setvalue{practical-chances.run.engine-curve.score}{0.4373}
\setvalue{practical-chances.run.gbm-context.score}{0.4260}
\setvalue{practical-chances.slice.clock.-30.engine-curve.said}{0.480}
\setvalue{practical-chances.slice.clock.-30.happened}{0.383}
\setvalue{practical-chances.slice.win.80-90.happened}{0.699}
\setvalue{puzzle-difficulty-unsolved.best.gain}{+45.01}
\setvalue{puzzle-difficulty-unsolved.best.score}{411.8}
\setvalue{puzzle-difficulty-unsolved.data.rows.test}{112{,}450}
\setvalue{puzzle-difficulty-unsolved.floor.score}{502.9}
\setvalue{puzzle-difficulty-unsolved.ladder.ahead.of}{all four networks}
\setvalue{puzzle-difficulty-unsolved.network.gain.of-lift}{98}
\setvalue{puzzle-difficulty-unsolved.reference.score}{456.8}
\setvalue{puzzle-difficulty-unsolved.shown.best.kept.percent}{38}
\setvalue{puzzle-difficulty-unsolved.shown.best.worth}{150}
\setvalue{puzzle-difficulty.best.gain}{+32.98}
\setvalue{puzzle-difficulty.best.mae}{196.8}
\setvalue{puzzle-difficulty.best.reliability.widest-gap}{61.3}
\setvalue{puzzle-difficulty.best.score}{261.9}
\setvalue{puzzle-difficulty.best.spearman}{0.8601}
\setvalue{puzzle-difficulty.context.lift.share}{36}
\setvalue{puzzle-difficulty.data.rows}{1{,}127{,}915}
\setvalue{puzzle-difficulty.data.rows.test}{112{,}450}
\setvalue{puzzle-difficulty.data.rows.train}{958{,}941}
\setvalue{puzzle-difficulty.floor.score}{502.9}
\setvalue{puzzle-difficulty.ladder.ahead.of}{all four networks}
\setvalue{puzzle-difficulty.ladder.count}{four}
\setvalue{puzzle-difficulty.recipe.sample.percent}{25}
\setvalue{puzzle-difficulty.reference.score}{294.9}
\setvalue{puzzle-difficulty.run.conv-small.parameters.millions}{1.0}
\setvalue{puzzle-difficulty.run.transformer-medium-long.epochs}{24}
\setvalue{puzzle-difficulty.run.transformer-medium.epochs}{12}
\setvalue{puzzle-difficulty.run.transformer-medium.parameters.millions}{6.6}
\setvalue{puzzle-difficulty.run.transformer-small.parameters.millions}{0.9}
\setvalue{puzzle-difficulty.task.metric}{RMSE}
\setvalue{tactic-found.best.bottom.10.rate}{67.2}
\setvalue{tactic-found.best.calibration.widest-gap}{1.1}
\setvalue{tactic-found.best.score}{0.5369}
\setvalue{tactic-found.best.top.10.rate}{95.6}
\setvalue{tactic-found.context.lift.share}{35}
\setvalue{tactic-found.data.rows}{810{,}140}
\setvalue{tactic-found.floor.constant}{69.0\%}
\setvalue{tactic-found.floor.score}{0.6203}
\setvalue{tactic-found.reference.score}{0.5414}
\setvalue{tactic-found.run.puzzle-rating.lift.share}{40}
\setvalue{tactic-found.slice.sacrifice.no.happened}{72.0\%}
\setvalue{tactic-found.slice.sacrifice.yes.happened}{51.0\%}
\setvalue{think-time.best.mae.factor}{1.86}
\setvalue{think-time.best.r2}{0.3144}
\setvalue{think-time.context.lift.share}{66}
\setvalue{think-time.floor.mae.factor}{2.16}
\setvalue{think-time.reference.mae.factor}{1.88}
\setvalue{think-time.run.gbm-context.leans.1.feature}{the mover's clock}
\setvalue{think-time.run.transformer-medium-long.gain.of-lift}{2}


\begin{document}

\papertitle{Predicting Human Behaviour in Chess: A Multi-Task Comparison of
  Gradient-Boosted Trees and Neural Networks}
  {Results as of \V{lab.results.date}\quad\textperiodcentered\quad
   \V{lab.side.count} tasks\quad\textperiodcentered\quad
   \V{lab.runs.count} models and \V{lab.studies.count} ablation models}

\begin{abstract}
Chess engines estimate the value of a position under best play. They do not
describe how people play it: how hard a tactic is to find, how likely a player
is to blunder, how long a move takes, or which of two rated players wins. We
study \V{lab.side.count} such prediction tasks on open Lichess data under a
single protocol. For every task we train a constant baseline, gradient-boosted
decision trees (GBDT) on scalar context features such as ratings and clocks,
GBDT with hand-crafted position features, and convolutional and transformer
networks that read the board as 64 square tokens. Every comparison is a paired
bootstrap over the same test games and is reported with a 95\% confidence
interval. The best neural network is significantly better than the GBDT
baseline in \V{lab.network.ahead.of}. On a normalised scale where the constant
baseline scores 0 and the GBDT baseline scores 100, its gain ranges from
\V{lab.network.gain.whole.least} to \V{lab.network.gain.whole.most} points on
the \V{lab.whole.count} tasks in which the models see every available input,
and reaches \V{lab.network.gain.most} points on the task that withholds the
puzzle solution, where hand-crafted features have little to work with. Context
features alone account for between \V{lab.context.share.least}\% and
\V{lab.context.share.most}\% of the GBDT baseline's improvement over the
constant baseline. Initialising the transformer encoder from a model
pretrained to predict engine evaluations is the most consistent improvement we
measure: it is significantly better in \V{lab.change.start.net.better.of} and
significantly worse in \V{lab.change.start.net.worse.of}. Larger training
sets improve both model classes. Each network was trained with a single random
seed, and the game data come from one month of rated rapid games on one site,
which limits how far the smaller differences generalise.
\end{abstract}

\section{Introduction}
\label{sec:intro}

Chess engines play far better than any person, and for any position they
return the best move and the expected result under best play. A coaching
application needs something else: a description of how people play. Whether a
player of a given rating will find a tactic, where such players usually
blunder, how long they spend on a move, and how a game between two rated
players is likely to end are properties of human behaviour. An engine
evaluation does not determine them.

Models of human move choice, notably Maia~\cite{maia} and
Maia-2~\cite{maia2}, address one of these questions. The others have received
less attention, and published results are hard to compare because each study
uses its own data, inputs and metrics. This paper reports one study of
\V{lab.side.count} supervised prediction tasks about human play. All of them
are built from the open Lichess database~\cite{lichessdb} and evaluated under
the same protocol. The work supports a chess coaching application that
reviews one player's games, so the questions are practical: for each task,
whether a learned model is accurate enough to be useful, how much of that
accuracy requires a neural network, and what would improve it.

Each task has a dataset split by game, a primary metric fixed before training,
and the same sequence of models: a constant baseline, gradient-boosted
decision trees (GBDT)~\cite{lightgbm} on scalar context features, GBDT with
hand-crafted position features, and residual convolutional and transformer
networks that read the board as square tokens. Because the sequence is fixed,
we can measure what each component contributes: the context features, the
hand-crafted features, the network, and then individual changes to the
network's training schedule, initialisation, architecture and training-set
size. Every comparison is paired on the same resampled test games.

Our main findings are the following.
\begin{itemize}
\item \textbf{The value of context varies widely between tasks.} GBDT on
ratings, clocks and other scalar context features, with no view of the board,
achieves between \V{lab.context.share.least}\% and
\V{lab.context.share.most}\% of the improvement that GBDT with hand-crafted
features achieves over the constant baseline (Section~\ref{sec:context}).
\item \textbf{Neural networks beat GBDT on most tasks, by a modest margin.}
The best network is significantly better than the GBDT baseline in
\V{lab.network.ahead.of}. Its gain is at most
\V{lab.network.gain.whole.most} points of normalised score where the models
see every available input, and \V{lab.network.gain.most} points where the
puzzle solution is withheld (Section~\ref{sec:networks}).
\item \textbf{Pretraining on engine evaluations is the most consistent
improvement.} A transformer whose encoder is initialised from a model
pretrained on \V{position-evaluation.data.rows.train} engine-evaluated
positions is significantly better than the same transformer trained from
scratch in \V{lab.change.start.net.better.of}. The exceptions are in
\V{lab.change.start.net.worse.models}, where the engine evaluation given as
an input is already close to the target (Section~\ref{sec:start}).
\item \textbf{Architecture matters less than initialisation.} An
attack-aware attention bias is significantly better in
\V{lab.change.attacks.net.better.of}, and a residual convolutional network of
similar size is significantly better than the transformer in
\V{lab.change.conv.net.better.of} while training faster. A
convolution-transformer hybrid and a larger transformer do not help
consistently (Section~\ref{sec:shapes}).
\item \textbf{More training data helps both model classes.} Halving the
training set costs the transformer more than GBDT in every task where both
were measured. Doubling and quadrupling it, at a fixed number of
optimisation steps, helps GBDT more than the transformer in
\V{lab.more.transformer-medium-hand.trees.of} (Section~\ref{sec:games}).
\item \textbf{The baselines are cheap.} The constant, formula and GBDT
models of all \V{lab.models.count} tasks train in
\V{lab.cost.seconds.others} on a CPU; the \V{lab.networks.count} networks
required \V{lab.cost.seconds.networks} on one GPU
(Section~\ref{sec:cost}).
\end{itemize}

Section~\ref{sec:each} summarises the result for each task,
Section~\ref{sec:use} discusses which models are accurate enough to use in an
application, and Section~\ref{sec:limits} lists the limitations of the study.
Most tasks have a separate paper with full results, ablations and an error
analysis; \V{lab.paperless.names} do not yet, and are documented by their
model cards.

\section{Related Work}
\label{sec:related}

\paragraph{Predicting human moves.}
Maia~\cite{maia} trains one neural network per rating band on Lichess games
and predicts human moves much more accurately than engines do.
Maia-2~\cite{maia2} replaces the per-band networks with a single model
conditioned on both players' ratings. Allie~\cite{allie} trains a transformer
on human games to predict moves, thinking time and resignation, and combines
it with a time-adaptive search. We include Maia-2, as released, as a baseline
in our human-move task. We do not expect to match it, because it was trained
on several orders of magnitude more games.

\paragraph{Human error and skill.}
Anderson et al.~\cite{humanerror} label errors exactly using endgame
tablebases and find that the difficulty of the position predicts errors
better than the player's skill or the time available. Regan and
Haworth~\cite{intrinsic} estimate a player's rating from the engine-judged
quality of the moves played. Our blunder-hazard, tactic-found and game-rating
tasks ask related questions about ordinary rated games, with ratings and
clocks as inputs.

\paragraph{Puzzle difficulty.}
Predicting the rating of Lichess puzzles has been the subject of a
data-mining competition~\cite{puzzlecup} and of a dedicated transformer
model~\cite{glickformer}. Both predict the puzzle's
Glicko-2 rating~\cite{glicko2} from its position and solution. We study that
setting and, in addition, the setting in which the solution is withheld.

\paragraph{Neural networks for chess positions.}
AlphaZero~\cite{alphazero} established residual convolutional networks for
evaluating chess positions, and Ruoss et al.~\cite{searchless} show that a
large transformer trained by supervised learning on engine evaluations plays
strongly without search. We train a far smaller model on engine evaluations,
and use it only as a source of pretrained weights for the other tasks.

\paragraph{Engine evaluations as win probabilities.}
Lichess converts an engine evaluation into a win probability with a logistic
curve fitted to game results~\cite{lichessaccuracy,lilawinchances}. We use
this curve to define a blunder and as a baseline for predicting game
outcomes.

\section{Tasks and Data}
\label{sec:models}

\subsection{Data sources}

All data come from the Lichess open database, which is released into the
public domain~\cite{lichessdb}: rated games with clock times and, for a
subset, the server's engine analysis; and the puzzles that Lichess extracts
from its games, each with a Glicko-2 rating~\cite{glicko2} derived from
solvers' attempts. One dataset joins the puzzles to the games they were
extracted from~\cite{puzzlegames}. No private games are used for training or
evaluation.

\subsection{Tasks}

Table~\ref{tab:models} lists the tasks. Each asks for a quantity that
describes human play and that an engine does not compute. Throughout, the
\emph{mover} is the side to move in a position.

\begin{table}[htbp]\centering\small
\block{lab-models}
\caption{The tasks: the prediction target, the number of examples in the
dataset, the number of examples in the test set and of games they come from,
and the primary metric ($\downarrow$: lower is better; $\uparrow$: higher is
better). The positions of \V{lab.bodies.names}, the pretraining task, are
not grouped into games, so each counts as its own game.}
\label{tab:models}
\end{table}

\begin{itemize}
\item \textbf{Puzzle difficulty.} Given a puzzle's position and its solution,
predict the puzzle's Lichess rating.
\item \textbf{Puzzle difficulty unsolved.} The same puzzles and target with
the solution withheld. The models see what the solver sees: the position and
the opponent's last move. Comparing the two tasks model by model measures how
much of a puzzle's difficulty can be predicted without knowing the solution.
\item \textbf{Blunder hazard.} Before a move is played, predict the
probability that the mover blunders, defined as losing more than
\V{method.blunder-bar} percentage points of win probability with the move.
Inputs are the position, both ratings, both clocks and the engine evaluation
before the move.
\item \textbf{Think time.} Predict how long the mover spends on the move,
given the position, the ratings and the clocks. No engine output is
given, since the aim is to model what a player can see.
\item \textbf{Tactic found.} In the game a puzzle was extracted from, predict
whether the player found the first move of the tactic.
\item \textbf{Practical chances.} Predict the expected game result for the
mover from the position, both ratings, both clocks and the engine
evaluation. The task is compared against the win probability implied by the
engine evaluation alone.
\item \textbf{Where to look.} Given a puzzle position and the opponent's last
move, but not the solution, predict for each of the
\V{where-to-look.task.candidates.count} squares whether a move of the
solution starts or ends on it.
\item \textbf{Puzzle themes.} Given a puzzle's position and solution, predict
which of \V{puzzle-themes.task.candidates.count} Lichess theme tags the
puzzle carries, such as fork, pin or deflection.
\item \textbf{Human move.} Given the position, both ratings and both clocks,
and no engine output, predict a probability for each legal move.
Maia-2~\cite{maia2} is evaluated on the same test set.
\item \textbf{Game rating.} Given one player's moves in a game, each with its
clock times and engine-judged loss, and neither player's rating, predict
that player's rating.
\end{itemize}

Blunder hazard, think time, practical chances and human move share one set of
positions: \V{blunder-hazard.data.rows} positions from
\V{blunder-hazard.recipe.games} rated \V{blunder-hazard.recipe.speeds} games
played in \V{blunder-hazard.data.month} that carry clock times and engine
analysis, with at most \V{blunder-hazard.recipe.per-game} positions drawn
pseudo-randomly from a game, from ply \V{blunder-hazard.recipe.min-ply} on.
The puzzle tasks share \V{puzzle-difficulty.data.rows} puzzles, a
\V{puzzle-difficulty.recipe.sample.percent}\% sample of the Lichess puzzles
whose rating has settled (a rating deviation of at most
\V{puzzle-difficulty.recipe.filters.max-rating-deviation} after at least
\V{puzzle-difficulty.recipe.filters.min-plays} attempts). Tactic found uses
\V{tactic-found.data.rows} puzzles under the same filter from a dataset that
joins each puzzle to its source game~\cite{puzzlegames}, restricted to
blitz, rapid and classical games. Game rating uses \V{game-rating.data.rows} examples
from \V{game-rating.recipe.games} rated \V{game-rating.recipe.speeds} games
of \V{game-rating.data.month}, one for each player with at least
\V{game-rating.recipe.least-moves} moves that carry a clock time and an
engine evaluation, for \V{game-rating.data.positions} positions in total.

\paragraph{The pretraining task.}
Table~\ref{tab:models} also lists \V{lab.bodies.names}: predict the win
probability that the engine assigns to a position. This quantity describes
the engine and not human play, so we exclude the task from every comparison
between tasks. We use it to pretrain a transformer encoder on
\V{position-evaluation.data.rows.train} positions from the Lichess file of
engine evaluations~\cite{lichessdb} (Section~\ref{sec:start}).

\subsection{Splits}

A dataset's examples are positions, and positions from the same game are
correlated, so we split by game: all positions of a game fall on the same
side of the split. For game rating both players of a game fall on the same
side. A hash of the game identifier assigns \V{lab.split.test.percent}\% of
games to the test set and \V{lab.split.valid.percent}\% to the validation
set. The pretraining dataset, whose positions are not grouped into games, is
split by position, with \V{lab.split.apart.test.percent}\% for testing and
\V{lab.split.apart.valid.percent}\% for validation. Model selection,
including the choice of training epoch, uses only the validation set. The
test set is used once, to report the final score.

\section{Methods}
\label{sec:harness}

All tasks share one implementation of the data pipeline, the input encoding,
the models, the training loop and the evaluation. A task is defined by a
configuration: its dataset, its target, the context features a model may use,
and its primary metric.

\subsection{Input representation}

A position is encoded as 64 tokens, one per square, each identifying the
piece on the square. The board is oriented so that the mover always
plays up the board. A task may add \emph{auxiliary input planes} that mark
squares, for example the squares of the opponent's last move or of the
puzzle's solution. \emph{Scalar context features} describe the situation
around the position: the two ratings, the two clocks and the move number,
and an engine evaluation in the tasks that allow one. \emph{Hand-crafted
features} are counts computed from the position: material, mobility, hanging
pieces, passed pawns and king safety, and for the puzzle tasks features of
the solution line and of its first move. For game rating the hand-crafted
features are aggregates over the player's moves: engine-judged move
accuracy, clock use and the number of available moves.

\subsection{Models}
\label{sec:sequence}

Every task is given the same sequence of models, so that each step measures
the contribution of one component.
\begin{itemize}
\item \run{constant}, the constant baseline, predicts the same value for
every example: the mean or median of the training targets, or the base rate.
\item \run{gbm-context} is GBDT (LightGBM~\cite{lightgbm}) on the scalar
context features only. It never sees the board.
\item \run{gbm-hand} is GBDT on the context features and the hand-crafted
features. It is the reference baseline, which we refer to as \emph{the GBDT
baseline}: every other model is compared against it.
\item \run{conv-small} and \run{transformer-small} are a residual
convolutional network (ResNet)~\cite{alphazero} and a
transformer~\cite{attention} over the square tokens, with
\V{puzzle-difficulty.run.conv-small.parameters.millions} and
\V{puzzle-difficulty.run.transformer-small.parameters.millions} million
parameters.
\item \run{transformer-medium} is the same transformer with
\V{puzzle-difficulty.run.transformer-medium.parameters.millions} million
parameters, and \run{transformer-medium-hand} additionally receives the
hand-crafted features.
\end{itemize}
The transformer uses a learned relative position bias, one per attention
head, for each offset between two squares~\cite{relative}. Where a formula
answers a task's question without any fitting, it is included as a further
baseline: the win probability implied by the engine evaluation, or the Elo
expected score computed from two ratings.

Three task types need a different output. Where the target is a set
(where to look, puzzle themes), a network outputs a probability for every
square or tag, and the constant baseline outputs each candidate's frequency
in the training set. For where to look, GBDT is trained with one example per
square, and for puzzle themes with one model per tag. For human move, a
network scores each legal move from its origin and destination squares, GBDT
is trained with one example per legal move, and the constant baseline is
uniform over the legal moves. Because one position yields many GBDT
examples, the GBDT models of where to look and human move are trained on a
subsample: \V{where-to-look.run.gbm-hand.train-rows} of
\V{where-to-look.data.rows.train} training puzzles and
\V{human-move.run.gbm-hand.train-rows} of \V{human-move.data.rows.train}
training positions. For game rating, the input is a sequence of positions:
the encoder embeds each position, and additional transformer layers operate
on the sequence of embeddings.

By default GBDT is trained with a learning rate of \V{method.gbm.rate} and
\V{method.gbm.leaves} leaves per tree for at most \V{method.gbm.rounds}
boosting rounds, and stops early when the validation score has not improved
for \V{method.gbm.patience} rounds; the tasks with structured outputs use a
higher learning rate and fewer rounds. Networks are trained with
AdamW~\cite{adamw} and a cosine learning-rate schedule with warm-up. After
each epoch the model is evaluated on the validation set, the weights of the
best epoch are kept, and training stops early when the validation score stops
improving. The paper of each task lists all hyperparameters.

\subsection{Modifications}
\label{sec:modifications}

Beyond the fixed sequence, we evaluate the following modifications, each
applied to one of the two medium transformers with nothing else changed
unless stated.
\begin{itemize}
\item \textbf{Longer schedule} (suffix \texttt{-long}): twice the training
epochs.
\item \textbf{Pretraining} (suffix \texttt{-pretrained}): the encoder is
initialised from the pretraining task instead of from random weights.
\item \textbf{Attack-aware attention bias} (suffix \texttt{-attacks}): each
attention head learns additional biases for a pair of squares when the piece
on one attacks the other in the current position, either directly or through
exactly one intervening piece, as in a pin or a skewer. The relations are
computed from the tokens without search.
\item \textbf{No relative position bias} (suffix \texttt{-no-geometry}): the
relative position bias is removed.
\item \textbf{Other encoders}: a ResNet of similar size with
squeeze-and-excitation (\run{conv-medium}), a hybrid of residual blocks
followed by transformer layers (\run{hybrid-medium}), and a larger
transformer (\run{transformer-large}).
\item \textbf{Training-set size}: half and a quarter of the training games,
and twice and four times as many from a larger sample of the same source.
\item \textbf{Feature-group ablation}: the hand-crafted features with one
group removed.
\end{itemize}

\section{Evaluation Protocol}
\label{sec:protocol}

\paragraph{Metrics.}
Each task has one primary metric, chosen before the first full training run
(Table~\ref{tab:models}): root-mean-square error (RMSE) or mean absolute error
(MAE) for the regression tasks, the area under the precision-recall curve
(PR-AUC) or the log loss for the binary tasks, average precision over the
candidates for the set-valued tasks, and top-1 accuracy for human move.

\paragraph{Calibration.}
For the binary and the set-valued tasks, the probabilities of every model are
rescaled by Platt scaling~\cite{platt} before they are scored: a logistic
function of the log-odds whose slope and intercept are fitted on the
validation set. For human move, a single softmax temperature is fitted in the
same way. Both transformations are monotone, so they change a model's log
loss and calibration and leave its ranking of examples unchanged.

\paragraph{Confidence intervals and paired comparisons.}
We report every score with a 95\% confidence interval (CI) obtained from
\V{method.resamples} bootstrap resamples of the test games. To compare two
models we compute the difference of their scores on each resample, which
gives a paired difference with its own interval. We call the difference a
\emph{gain} and sign it so that a positive gain means the first model is
better. A gain is \emph{statistically significant} when its 95\% CI excludes
zero. These intervals reflect sampling of the test games only. They do not
include the variation between training runs with different random seeds.

\paragraph{Acceptance criterion.}
A modification is \emph{accepted} only if the modified model is better than
the unmodified one on the validation set and its paired gain on the test set
is significantly greater than zero. A modification that is not accepted is
still reported, as a measurement of its effect.

\paragraph{Normalised score.}
The tasks use different metrics, so to compare across tasks we normalise a
model's score $s(m)$ by the two fixed baselines:
\[
  N(m) = 100 \times \frac{s(m) - s(\text{constant})}{s(\text{GBDT baseline}) - s(\text{constant})}.
\]
The constant baseline scores 0 and the GBDT baseline scores 100 on every
task, and a difference between two models is expressed in \emph{points} of
normalised score. An interval on this scale is the interval of the paired
gain divided by the same denominator. The denominator is itself an estimate
and we do not propagate its uncertainty, so the normalised score is used
only for comparison across tasks. The paper of each task reports all gains in
the task's own units.

\section{Results}
\label{sec:across}

\subsection{Main results}

Table~\ref{tab:scores} gives the main results for each task in its own
metric, and Figure~\ref{fig:shares} places the models of the fixed sequence
on the normalised scale.

\begin{table}[htbp]\centering\small
\block{lab-scores}
\caption{Main results on the test set: the constant baseline, GBDT on context
features, GBDT with hand-crafted features (the GBDT baseline), and the best
neural network with its paired gain over the GBDT baseline, in the task's
metric (\emph{Gain}) and in points of normalised score (\emph{Norm.\ gain}).
A gain in bold is statistically significant.
Figure~\ref{fig:networks} shows the confidence intervals, and
Appendix~\ref{sec:runs} lists every model.}
\label{tab:scores}
\end{table}

\begin{figure}[htbp]\centering
\block{lab-shares}
\caption{Normalised score of each model of the fixed sequence and of each
formula baseline, by task. Zero is the constant baseline and 100 is GBDT with
hand-crafted features. Each model keeps its colour and marker across tasks.}
\label{fig:shares}
\end{figure}

\subsection{Context features}
\label{sec:context}

The contribution of the scalar context features varies more between tasks
than any other quantity in Figure~\ref{fig:shares}. GBDT on context features
alone reaches a normalised score of
\V{puzzle-difficulty.context.lift.share} in puzzle difficulty,
\V{tactic-found.context.lift.share} in tactic found,
\V{blunder-hazard.context.lift.share} in blunder hazard,
\V{think-time.context.lift.share} in think time and
\V{practical-chances.context.lift.share} in practical chances. In the tasks
with structured outputs, the context-only GBDT also receives candidate-level
inputs that need no chess knowledge, such as the location of a square and the
piece on it, or the piece that moves and its origin and destination. It
reaches \V{where-to-look.run.gbm-context.lift.share} in where to look and
\V{human-move.run.gbm-context.lift.share} in human move. In puzzle themes,
where the context is the length of the solution and the castling rights, it
reaches \V{puzzle-themes.run.gbm-context.lift.share}, and in game rating,
where it is the time control and the length of the game,
\V{game-rating.run.gbm-context.lift.share}.

We hypothesise that this ordering reflects what determines each target. The
difficulty of a tactic depends mostly on the position and the solution, and
little on the solver. Thinking time and blunder rate depend strongly on the
situation (the ratings, the clocks and, for blunders, the current
evaluation), whatever the arrangement of the pieces. In practical chances one
of the
context features is the engine evaluation, which is itself an estimate of the
target: used directly as a prediction with no fitting, it reaches a
normalised score of \V{practical-chances.run.engine-curve.lift.share}. For an
application, a high score for the context-only model means that a lookup
table over the context features captures most of what any model learns on
that task.

\subsection{Neural networks compared with GBDT}
\label{sec:networks}

Figure~\ref{fig:networks} compares every network with the GBDT baseline.

\begin{figure}[htbp]\centering
\block{lab-networks}
\caption{Paired gain of each network over the GBDT baseline with its 95\%
CI, in points of normalised score. Each panel is a task and has its own
horizontal scale. A network to the right of the vertical line is better than
the GBDT baseline, and the gain is significant when the interval does not
cross the line.}
\label{fig:networks}
\end{figure}

The best network is significantly better than the GBDT baseline in
\V{lab.network.ahead.of} (\V{lab.network.ahead.names}), not significantly
different in \V{lab.network.level.of} (\V{lab.network.level.names}) and
significantly worse in \V{lab.network.behind.of}. Among the
\V{lab.whole.count} tasks in which the models see every available input, the
largest gain is \V{lab.network.gain.whole.most} points, in
\V{lab.network.gain.whole.most.model}. In these tasks GBDT with hand-crafted
features therefore already achieves most of the improvement over the
constant baseline, and the network adds a smaller amount. The exception is
\V{lab.hidden.names}, where the solution is withheld from all models. There
the hand-crafted features describe only the position, the GBDT baseline is
weak, and the best network gains
\V{puzzle-difficulty-unsolved.network.gain.of-lift} points. This number
reflects the weakness of that baseline as much as the strength of the
network.

In \V{lab.network.ahead.run.transformer-medium-hand-pretrained.names} the
best network is the medium transformer with hand-crafted features and a
pretrained encoder, and in
\V{lab.network.ahead.run.transformer-medium-attacks.names} it is the
transformer with the attack-aware attention bias. In
\V{lab.network.level.run.transformer-medium-hand-quadruple.names} it is the
transformer with hand-crafted features trained on four times the standard
training set, and in the remaining tasks it is a network of the fixed
sequence (Table~\ref{tab:scores}). The best model overall is a neural network
in \V{lab.best.network.of}. The exceptions are human move, where Maia-2
leads, and practical chances, where GBDT trained on a larger training set
leads.

\paragraph{Networks without hand-crafted features.}
Whether a network that reads only the board and the context features can
beat GBDT with hand-crafted features depends on its initialisation. Trained
from scratch, \run{transformer-medium} is significantly better than the GBDT
baseline in \V{lab.ladder.board.ahead.of} (\V{lab.ladder.board.ahead.names}),
not significantly different in \V{lab.ladder.board.level.of}
(\V{lab.ladder.board.level.names}) and significantly worse in
\V{lab.ladder.board.behind.of} (\V{lab.ladder.board.behind.names}). The two
small networks perform similarly: the better of them is significantly better
than the GBDT baseline in \V{lab.small.ahead.of}. Counting every
modification, the best network without hand-crafted features is
significantly better than the GBDT baseline in \V{lab.board.ahead.of} and
significantly worse in \V{lab.board.behind.of}. Where pretraining was
accepted, the pretrained transformer without hand-crafted features is
significantly better than the GBDT baseline in
\V{lab.started.board.ahead.of} (\V{lab.started.board.ahead.names}),
compared with \V{lab.started.ladder.board.ahead.of} for the same transformer
trained from scratch.

Trained from scratch, the network without hand-crafted features beats GBDT
mainly on the puzzle tasks and on the tasks with structured outputs. Two
explanations are consistent with this. First, the puzzle datasets are the
largest (Table~\ref{tab:models}), so a network has the most positions to
learn from. Second, a tactic is a spatial pattern of pieces, which a network
can read directly, whereas blunders and thinking time depend more on
whole-board quantities such as material balance and hanging pieces, which
the hand-crafted features provide to GBDT directly. Human move separates the
two explanations to some degree: it uses the same positions as blunder
hazard and think time, with no additional data, and there the transformer
trained from scratch is significantly better than GBDT. The GBDT baseline of
human move is trained on \V{human-move.run.gbm-hand.train-rows} of the
\V{human-move.data.rows.train} training positions
(Section~\ref{sec:sequence}), so the fairer comparison is with the transformer
trained on a quarter of the games, which also remains ahead (gain
\V{human-move.run.transformer-medium-quarter.gain}, 95\% CI
\V{human-move.run.transformer-medium-quarter.gain.low} to
\V{human-move.run.transformer-medium-quarter.gain.high}). The data-size
ablation of Section~\ref{sec:games} supports the first explanation as well.

\subsection{Hand-crafted features}
\label{sec:features}

Table~\ref{tab:changes} and Figure~\ref{fig:changes} report three
modifications as paired gains across tasks: adding the hand-crafted features,
doubling the training epochs, and pretraining the encoder.

\begin{table}[htbp]\centering\small
\block{lab-changes}
\caption{Paired gain of a modified model over the unmodified model named in
the column head, by task, in points of normalised score with the 95\% CI.
Gains in bold are statistically significant. The paper of each task reports
the same gains in the task's own metric.}
\label{tab:changes}
\end{table}

\begin{figure}[htbp]\centering
\block{lab-changes-figure}
\caption{The paired gains of Table~\ref{tab:changes} with their 95\% CIs.}
\label{fig:changes}
\end{figure}

For GBDT, adding the hand-crafted features to the context features is
\V{lab.change.hand.over.gbm-context.said}. For the medium transformer trained
from scratch, adding them is \V{lab.change.hand.over.transformer-medium.said}.
With the longer schedule the same comparison is
\V{lab.change.hand.over.transformer-medium-long.said}. In game rating the
features make the transformer worse, which we discuss in
Section~\ref{sec:game-rating}.

To see which features matter, we retrained the GBDT baseline with one group
of hand-crafted features removed at a time (Table~\ref{tab:features}). A
significant loss shows that GBDT used the group and could not recover the
same information from its other inputs. No loss means either that the group
is not useful or that it is redundant with other features. The group whose
removal costs the most is \V{lab.features.in.puzzle-difficulty.most.group}
in puzzle difficulty (\V{lab.features.in.puzzle-difficulty.most.gain}
points), \V{lab.features.in.tactic-found.most.group} in tactic found
(\V{lab.features.in.tactic-found.most.gain}),
\V{lab.features.in.blunder-hazard.most.group} in blunder hazard
(\V{lab.features.in.blunder-hazard.most.gain}),
\V{lab.features.in.think-time.most.group} in think time
(\V{lab.features.in.think-time.most.gain}),
\V{lab.features.in.practical-chances.most.group} in practical chances
(\V{lab.features.in.practical-chances.most.gain}) and
\V{lab.features.in.game-rating.most.group} in game rating
(\V{lab.features.in.game-rating.most.gain}).

\begin{table}[htbp]\centering\small
\block{lab-features}
\caption{Feature-group ablation of the GBDT baseline. Each cell is the paired
gain of the model without the group over the model with it, in points of
normalised score: a negative value means the model is worse without the
group, and bold marks a statistically significant difference. An empty cell
means the task's features do not include the group.}
\label{tab:features}
\end{table}

Removing material is \V{lab.features.material.said}, and removing passed
pawns is \V{lab.features.pawns.said}. Otherwise the tasks differ according
to their inputs. The two tasks that include a solution rely on its features:
the solution line as a whole in puzzle difficulty and its first move in
tactic found. The three tasks on game positions rely most on material, which
also indicates the phase of the game. In game rating the engine-judged
accuracy of the moves is by far the most important group.

\subsection{Training schedule}
\label{sec:epochs}

In the fixed sequence, \V{lab.ladder.at-last-epoch.of} reached their best
validation score in the final epoch, which suggests that some had not
converged. In \V{lab.change.epochs.over.transformer-medium.count} tasks we
therefore retrained both medium transformers for
\V{puzzle-difficulty.run.transformer-medium-long.epochs} epochs instead of
\V{puzzle-difficulty.run.transformer-medium.epochs}
(Table~\ref{tab:changes}). With the longer schedule, the transformer without
hand-crafted features is \V{lab.change.epochs.over.transformer-medium.said},
and the transformer with them is
\V{lab.change.epochs.over.transformer-medium-hand.said}.

The longer schedule helps most in puzzle difficulty, where it gains
\V{lab.change.epochs.over.transformer-medium.gain.most} points without the
hand-crafted features and
\V{lab.change.epochs.over.transformer-medium-hand.gain.most} points with
them. Elsewhere the gains are a few points or not significant. By the
acceptance criterion, \V{lab.change.epochs.net.kept.of} are accepted; both
models are accepted in \V{lab.change.epochs.net.kept.whole.models}.

\subsection{Pretraining on engine evaluations}
\label{sec:start}

The networks of the fixed sequence are trained from scratch, so everything
they know about how pieces move is learned from the task's own training set.
As an alternative we initialise the encoder from a model trained on a
different task. The pretraining task, \V{lab.bodies.names}, trains the
medium transformer without hand-crafted features to predict the engine's
win probability for a position, from evaluations at depth
\V{position-evaluation.recipe.min-depth} or more. On
\V{position-evaluation.data.rows.test} held-out positions its RMSE is
\V{position-evaluation.best.score} percentage points of win probability,
compared with \V{position-evaluation.reference.score} for GBDT with
hand-crafted features and \V{position-evaluation.floor.score} for the
constant baseline. Pretraining took \V{position-evaluation.best.seconds} on
one GPU and was done once. A pretrained model copies the piece and square
embeddings and all transformer blocks. The prediction head is initialised
randomly, and so is the projection of the context features, except in puzzle
difficulty unsolved, whose context features are those of the pretraining
task (the castling rights). The model is then trained on its task with the
standard schedule.

With a pretrained encoder, the transformer without hand-crafted features is
\V{lab.change.start.over.transformer-medium.said}, and the transformer with
them is \V{lab.change.start.over.transformer-medium-hand.said}. By the
acceptance criterion, \V{lab.change.start.net.kept.of} are accepted: both
models in \V{lab.change.start.net.kept.whole.models}. In each of these tasks
the pretrained transformer with hand-crafted features is the best model
(Table~\ref{tab:scores}). Four observations follow from
Table~\ref{tab:changes}.

\paragraph{Pretraining gains more than a longer schedule.}
In each of the \V{lab.change.epochs.over.transformer-medium.count} tasks that
have both modifications, both transformers gain more from pretraining than
from twice the epochs, and the confidence intervals of the two gains do not
overlap. Pretraining achieves this with half the training epochs.

\paragraph{The gain is largest without hand-crafted features.}
Wherever pretraining helps, it helps the transformer without hand-crafted
features more than the transformer with them. We hypothesise that part of
what pretraining provides overlaps with the hand-crafted features: an engine
evaluation depends on material and on what is attacked, so the pretrained
encoder has had to learn both from the board.

\paragraph{Pretraining does not help in practical chances.}
In practical chances both pretrained models are significantly worse, and
both reached their best validation epoch earlier than the models trained
from scratch (epochs
\V{practical-chances.run.transformer-medium-pretrained.best-epoch} and
\V{practical-chances.run.transformer-medium-hand-pretrained.best-epoch} of
\V{practical-chances.run.transformer-medium-pretrained.epochs}, compared with
\V{practical-chances.run.transformer-medium.best-epoch} and
\V{practical-chances.run.transformer-medium-hand.best-epoch}). The engine
evaluation is an input in this task, as it is in blunder hazard and tactic
found, where pretraining helps. What sets practical chances apart is that the
evaluation is close to the target itself: used directly as a prediction it
reaches a normalised score of
\V{practical-chances.run.engine-curve.lift.share}
(Section~\ref{sec:context}). We hypothesise that an encoder trained to
reproduce the evaluation therefore adds little that the model does not
already receive, whereas a blunder or a missed tactic depends on properties
of the position that the evaluation does not summarise.

\paragraph{Pretraining helps most where the most information is withheld.}
With the puzzle solution withheld, the paired gains from pretraining are
\V{puzzle-difficulty-unsolved.pair.transformer-medium-pretrained.transformer-medium.gain}
and
\V{puzzle-difficulty-unsolved.pair.transformer-medium-hand-pretrained.transformer-medium-hand.gain}
rating points of RMSE for the two transformers. With the solution given,
they are
\V{puzzle-difficulty.pair.transformer-medium-pretrained.transformer-medium.gain}
and
\V{puzzle-difficulty.pair.transformer-medium-hand-pretrained.transformer-medium-hand.gain}
on the same test puzzles.

We evaluated a single pretrained encoder, of one size and trained for
\V{position-evaluation.best.epochs} epochs on a sample of the available
evaluations. The results therefore measure one pretraining configuration and
not the best achievable one.

\subsection{Architecture}
\label{sec:shapes}

The architecture modifications of Section~\ref{sec:modifications} change the
encoder of a medium transformer and keep the training set, the input
features, the prediction head and the number of epochs. Each is compared
with the transformer it modifies (Table~\ref{tab:shapes},
Figure~\ref{fig:shapes}). All models in this section are trained from
scratch.

\begin{table}[htbp]\centering\small
\block{lab-shapes}
\caption{Architecture modifications: paired gain of the modified model over
the transformer named in the column head, by task, in points of normalised
score with the 95\% CI. Gains in bold are statistically significant.}
\label{tab:shapes}
\end{table}

\begin{figure}[htbp]\centering
\block{lab-shapes-figure}
\caption{The paired gains of Table~\ref{tab:shapes} with their 95\% CIs.}
\label{fig:shapes}
\end{figure}

\paragraph{Attack-aware attention bias.}
The relative position bias tells an attention head that two squares share a
diagonal, but not whether a piece stands between them. The attack-aware bias
adds this information. Without hand-crafted features, the transformer with
the bias is \V{lab.change.attacks.over.transformer-medium.said}. With
hand-crafted features it is
\V{lab.change.attacks.over.transformer-medium-hand.said}. By the acceptance
criterion, \V{lab.change.attacks.net.kept.of} are accepted.

In puzzle difficulty the bias gives the transformer without hand-crafted
features a gain of
\V{puzzle-difficulty.pair.transformer-medium-attacks.transformer-medium.gain}
rating points of RMSE, compared with
\V{puzzle-difficulty.pair.transformer-medium-pretrained.transformer-medium.gain}
from pretraining and
\V{puzzle-difficulty.pair.transformer-medium-long.transformer-medium.gain}
from the longer schedule. With the bias, the hand-crafted features add little
there: the two transformers with the bias score
\V{puzzle-difficulty.run.transformer-medium-hand-attacks.score} and
\V{puzzle-difficulty.run.transformer-medium-attacks.score}, each within the
other's confidence interval. In think time the bias moves the transformer
without hand-crafted features from behind the GBDT baseline to ahead of it:
its normalised gain over that baseline changes from
\V{think-time.run.transformer-medium.gain.of-lift} to
\V{think-time.run.transformer-medium-attacks.gain.of-lift} points. In
\V{lab.change.attacks.over.transformer-medium.level.models} the bias makes
no significant difference with or without hand-crafted features. These are
the two tasks with a binary target about a single move in a game, and also
the two tasks in which the longer schedule did not help the transformer with
hand-crafted features. A plausible explanation is that the weights of the
additional bias start at zero and must be learned from the task's own data,
and that these two tasks provide the weakest training signal for it.

\paragraph{ResNet.}
A ResNet of similar size to the medium transformer, with
squeeze-and-excitation, is trained at a higher peak learning rate, which is
the only other change. With hand-crafted features it is
\V{lab.change.conv.over.transformer-medium-hand.said}, and
\V{lab.change.conv.net.kept.of} are accepted. It also trains faster:
\V{puzzle-difficulty.run.conv-medium-hand.seconds} compared with
\V{puzzle-difficulty.run.transformer-medium-hand.seconds} in puzzle
difficulty, \V{think-time.run.conv-medium-hand.seconds} compared with
\V{think-time.run.transformer-medium-hand.seconds} in think time and
\V{tactic-found.run.conv-medium-hand.seconds} compared with
\V{tactic-found.run.transformer-medium-hand.seconds} in tactic found.
Because the learning rate changes together with the encoder, the gain cannot
be attributed to the convolutions alone. In human move, where the output is
a move and the models are trained without hand-crafted features, the ResNet
is \V{human-move.pair.conv-medium.transformer-medium.verdict} (gain
\V{human-move.pair.conv-medium.transformer-medium.gain} in top-1 accuracy).

\paragraph{Hybrid and larger transformer.}
With hand-crafted features, the hybrid encoder is
\V{lab.change.hybrid.over.transformer-medium-hand.said}. In human move,
without them, it is
\V{human-move.pair.hybrid-medium.transformer-medium.verdict}. By the
acceptance criterion, \V{lab.change.hybrid.net.kept.of} are accepted.

A transformer with
\V{puzzle-difficulty.run.transformer-large-hand.parameters.millions} million
parameters instead of
\V{puzzle-difficulty.run.transformer-medium-hand.parameters.millions}
million is
\V{puzzle-difficulty.pair.transformer-large-hand.transformer-medium-hand.verdict}
in puzzle difficulty, where it is trained with hand-crafted features (gain
\V{puzzle-difficulty.pair.transformer-large-hand.transformer-medium-hand.gain}
rating points of RMSE), and
\V{human-move.pair.transformer-large.transformer-medium.verdict} in human
move, where it is trained without them (gain
\V{human-move.pair.transformer-large.transformer-medium.gain} in top-1
accuracy). The larger model was trained with the learning rate and schedule
of the smaller one, so we cannot say whether a schedule of its own would
change these results.

\paragraph{Relative position bias.}
Removing the relative position bias from the transformer with hand-crafted
features is \V{lab.change.geometry.over.transformer-medium-hand.said}. In
human move, without hand-crafted features, removing it is
\V{human-move.pair.transformer-medium-no-geometry.transformer-medium.verdict}
by a wide margin (Section~\ref{sec:human-move}). The bias therefore matters
most where the output is a move.

\paragraph{Summary.}
No architecture modification is significant in \V{lab.shapes.still.names},
and a model with a modified architecture is the best model in
\V{lab.shapes.best.of} that have one. Where pretraining was also evaluated,
the pretrained transformer remains the best model. In human move, where it
was not, the transformer with the attack-aware bias is the best model
trained here, behind Maia-2. We have not combined an architecture
modification with pretraining.

\subsection{Training-set size}
\label{sec:games}

We retrained the GBDT baseline and the medium transformer with hand-crafted
features on half and on a quarter of the training games, and evaluated them
on the same test set (Table~\ref{tab:games}, Figure~\ref{fig:games}). The
subsets are nested. The transformer is given twice the epochs on half of the
games and four times the epochs on a quarter, so that the number of
optimisation steps stays constant and only the amount of data changes. We
also trained both models on twice and four times as many games, drawn from a
larger sample of the same source with the same split, so that the test set
is unchanged and no test game appears in the larger training sets
(Figure~\ref{fig:games-more}). Here the transformer is given proportionally
fewer epochs, again keeping the number of steps constant. Tactic found has
only the doubled training set, because its standard training set already
uses half of the available source.

\begin{table}[htbp]\centering\small
\block{lab-games}
\caption{Training-set size: paired gain of a model trained on a quarter,
half, twice or four times the standard number of training games over the
same model trained on the standard training set, in points of normalised
score. A negative value means the model is worse than with the standard
training set, and bold marks a statistically significant difference.}
\label{tab:games}
\end{table}

\begin{figure}[htbp]\centering
\block{lab-games-figure}
\caption{The paired gains of Table~\ref{tab:games} for smaller training
sets, with their 95\% CIs.}
\label{fig:games}
\end{figure}

\begin{figure}[htbp]\centering
\block{lab-games-more-figure}
\caption{The paired gains of Table~\ref{tab:games} for larger training
sets, with their 95\% CIs.}
\label{fig:games-more}
\end{figure}

\paragraph{Smaller training sets.}
On half of the training games the GBDT baseline is
\V{lab.games.gbm-hand.half.said}, and on a quarter it is
\V{lab.games.gbm-hand.quarter.said}. The transformer is significantly worse
in \V{lab.games.transformer-medium-hand.half.worse.of} on half of the games
and in \V{lab.games.transformer-medium-hand.quarter.worse.of} on a quarter
(\V{lab.games.transformer-medium-hand.quarter.worse.models}). For the same
reduction, the transformer loses more than GBDT in every task that has both
measurements. Extrapolated to larger training sets, this suggests that
additional data should narrow the gap between GBDT and the network, or widen
the network's lead.

\paragraph{Larger training sets.}
The larger training sets confirm that more data helps, but not that it
favours the network. On twice the games the GBDT baseline is
\V{lab.games.gbm-hand.double.said}, and on four times the games it is
\V{lab.games.gbm-hand.quadruple.said}. The transformer on twice the games is
\V{lab.games.transformer-medium-hand.double.said}, and on four times the
games it is significantly better in
\V{lab.games.transformer-medium-hand.quadruple.better.of}. Every model
trained on four times the games is accepted, and so is every model trained
on twice the games except the two whose gains are not significant. However,
the additional data helps GBDT more than the transformer in
\V{lab.more.transformer-medium-hand.trees.of}
(\V{lab.more.transformer-medium-hand.trees.models}). Only in
\V{lab.more.transformer-medium-hand.network.models} is the transformer's gain
the larger one.

\paragraph{Interpretation.}
We hypothesise that the extrapolation from smaller training sets failed
because the comparison holds the number of optimisation steps fixed. On half
and on a quarter of the games the transformer sees each game two and four
times as often, and in every task its best validation epoch came before the
end of its schedule, which indicates overfitting. On twice and four times
the games its best epoch was the last one in all but one model (twice the
games in practical chances), which indicates that it had not converged. Part of the transformer's loss on
smaller training sets is therefore due to overfitting under repeated passes
over the same games. That effect does not extend beyond the standard
training set, where the transformer is instead limited by the number of
steps. We have not yet trained a model with both more data and more steps.

Two further caveats apply. Each transformer was trained once, and its gains
on four times the games are small, between
\V{lab.games.transformer-medium-hand.quadruple.gain.least} and
\V{lab.games.transformer-medium-hand.quadruple.gain.most} points, so they
should be confirmed with additional seeds. And in puzzle difficulty the GBDT
model on four times the data built
\V{puzzle-difficulty.run.gbm-hand-quadruple.trees} trees, the maximum of
\V{method.gbm.rounds} boosting rounds we allow, so its gain may be limited
by that setting.

\subsection{Computational cost}
\label{sec:cost}

\begin{table}[htbp]\centering\small
\block{lab-cost}
\caption{Number of models and total training time by task. The constant,
formula and GBDT models were trained on a CPU (\V{lab.cost.cpus}) and each
network on one GPU (\V{lab.cost.gpus}). Ablation models are those trained
only to measure the effect of a modification.}
\label{tab:cost}
\end{table}

The constant, formula and GBDT models of all \V{lab.models.count} tasks
trained in \V{lab.cost.seconds.others} in total, and the
\V{lab.networks.count} networks in \V{lab.cost.seconds.networks}
(Table~\ref{tab:cost}). The \V{lab.studies.count} ablation models took a
further \V{lab.cost.seconds.studies}. A GBDT model is a small file that runs
on any device, whereas a network needs its weights and suitable hardware.
Where the models see every input, a network gains at most
\V{lab.network.gain.whole.most} points of normalised score over GBDT, which
is the trade-off to weigh when deploying a model.

\section{Results by Task}
\label{sec:each}

This section summarises the result of each task. The separate task papers
give every model, the full ablations and an error analysis with example
positions.

\subsection{Puzzle difficulty}

On \V{puzzle-difficulty.data.rows.test} test puzzles the best model has an
RMSE of \V{puzzle-difficulty.best.score} rating points, compared with
\V{puzzle-difficulty.reference.score} for the GBDT baseline and
\V{puzzle-difficulty.floor.score} for the constant baseline. Its Spearman
rank correlation with the true rating is \V{puzzle-difficulty.best.spearman}
and its MAE is \V{puzzle-difficulty.best.mae} points. The predictions are
well calibrated on average, within
\V{puzzle-difficulty.best.reliability.widest-gap} points across the rating
range, but the error on an individual puzzle is large
(Figure~\ref{fig:puzzle}). The model is therefore suitable for ranking the
tactics of a game by difficulty and not yet for reporting a rating for a
single tactic.

\begin{figure}[htbp]\centering
\block{puzzle-difficulty-reliability}
\caption{Puzzle difficulty: mean observed rating against mean predicted
rating, in equal-sized bins of the test set ordered by prediction. The
shaded band covers the 10th to 90th percentile of observed ratings in each
bin of the first model.}
\label{fig:puzzle}
\end{figure}

\subsection{Puzzle difficulty unsolved}

With the solution withheld and the same
\V{puzzle-difficulty-unsolved.data.rows.test} test puzzles, the best model
has an RMSE of \V{puzzle-difficulty-unsolved.best.score} rating points,
compared with \V{puzzle-difficulty-unsolved.reference.score} for the GBDT
baseline and \V{puzzle-difficulty-unsolved.floor.score} for the constant
baseline. Table~\ref{tab:unsolved} compares each model with its counterpart
that sees the solution. Between the best models of the two tasks, the
solution is worth \V{puzzle-difficulty-unsolved.shown.best.worth} rating
points of RMSE, and the model without the solution retains
\V{puzzle-difficulty-unsolved.shown.best.kept.percent}\% of the improvement
over the constant baseline.

\begin{table}[htbp]\centering\small
\block{puzzle-difficulty-unsolved-shown}
\caption{Puzzle difficulty with the solution withheld (\emph{Hidden}) and
given (\emph{Shown}): each model of the unsolved task beside the model with
the same configuration in puzzle difficulty, on the same test puzzles.
\emph{Difference} is the RMSE advantage of the model that sees the solution.
\emph{Retained} is the percentage of that model's improvement over the
constant baseline that the model without the solution achieves.}
\label{tab:unsolved}
\end{table}

Trained from scratch, no model retains half of the improvement of its
counterpart, and GBDT retains the least. Pretraining changes this: the two
pretrained transformers retain
\V{puzzle-difficulty-unsolved.shown.transformer-medium-hand-pretrained.kept.percent}\%
and
\V{puzzle-difficulty-unsolved.shown.transformer-medium-pretrained.kept.percent}\%,
compared with
\V{puzzle-difficulty-unsolved.shown.transformer-medium-hand.kept.percent}\%
and
\V{puzzle-difficulty-unsolved.shown.transformer-medium.kept.percent}\% for
the same models trained from scratch. How much of a puzzle's difficulty can
be predicted from the position alone thus depends on the model, and the
values reported here are lower bounds. This is also the task where a network
matters most: the gain of the best network over the GBDT baseline is
\V{puzzle-difficulty-unsolved.best.gain} rating points, compared with
\V{puzzle-difficulty.best.gain} when the solution is given.

\subsection{Blunder hazard}

Blunders are rare: the base rate is \V{blunder-hazard.floor.constant}. The
best model has a PR-AUC of \V{blunder-hazard.best.score}, compared with
\V{blunder-hazard.reference.score} for the GBDT baseline and
\V{blunder-hazard.floor.score} for the constant baseline. After calibration
the predicted probabilities match observed frequencies: none of the
\V{blunder-hazard.best.calibration.groups} bins deviates by more than
\V{blunder-hazard.best.calibration.widest-gap} percentage points
(Figure~\ref{fig:blunder}). Among the 1\% of test positions with the highest
predicted probability, \V{blunder-hazard.best.top.1.rate}\% were blunders,
\V{blunder-hazard.best.top.1.lift} times the base rate.

\begin{figure}[htbp]\centering
\block{blunder-hazard-calibration}
\caption{Blunder hazard: reliability diagram. Observed blunder frequency
against predicted probability in \V{blunder-hazard.best.calibration.groups}
equal-sized bins of the test set. A calibrated model lies on the diagonal.}
\label{fig:blunder}
\end{figure}

\subsection{Think time}

The target is the logarithm of the seconds spent, so an absolute error
corresponds to a multiplicative factor. The best model is off by a factor of
\V{think-time.best.mae.factor} on average, the GBDT baseline by
\V{think-time.reference.mae.factor} and the constant baseline by
\V{think-time.floor.mae.factor}. GBDT on context features alone reaches a
normalised score of \V{think-time.context.lift.share}, and its most
important feature is \V{think-time.run.gbm-context.leans.1.feature}
(Figure~\ref{fig:think}). Most of the variation in thinking time is not
predictable from these inputs: the best model has an $R^2$ of
\V{think-time.best.r2}. The model describes what is typical, so a single
move that deviates from it is weak evidence, whereas a consistent pattern
over many moves is informative.

\begin{figure}[htbp]\centering
\block{think-time-slice-clock-said}
\caption{Think time: typical time spent on a move as a function of the time
remaining on the mover's clock, with the mean predictions of two
models.}
\label{fig:think}
\end{figure}

\subsection{Tactic found}

In the source games the tactic was found in
\V{tactic-found.floor.constant} of cases. The best model has a log loss of
\V{tactic-found.best.score}, compared with \V{tactic-found.reference.score}
for the GBDT baseline and \V{tactic-found.floor.score} for the constant
baseline. The Elo expected score computed from the puzzle's rating and the
player's rating reaches a normalised score of only
\V{tactic-found.run.puzzle-rating.lift.share}. The outcome depends more on
the tactic than on the player: in Figure~\ref{fig:tactic} the share found
varies far more with the puzzle's rating than with the player's rating. A
first move that sacrifices material was found in
\V{tactic-found.slice.sacrifice.yes.happened} of cases, compared with
\V{tactic-found.slice.sacrifice.no.happened} otherwise. Among the 10\% of
tactics with the highest predicted probability,
\V{tactic-found.best.top.10.rate}\% were found; among the 10\% with the
lowest, \V{tactic-found.best.bottom.10.rate}\% were missed.

\begin{figure}[htbp]\centering
\block{tactic-found-grid-ratings}
\caption{Tactic found: percentage of tactics found in the game, by the
rating of the mover and the rating of the puzzle.}
\label{fig:tactic}
\end{figure}

\subsection{Practical chances}

Used directly as a prediction of the game result, the win probability
implied by the engine evaluation has an RMSE of
\V{practical-chances.run.engine-curve.score}. GBDT on that evaluation
together with the ratings, the clocks and the move number reaches
\V{practical-chances.run.gbm-context.score}, a paired gain of
\V{practical-chances.pair.gbm-context.engine-curve.gain} (95\% CI
\V{practical-chances.pair.gbm-context.engine-curve.gain.low} to
\V{practical-chances.pair.gbm-context.engine-curve.gain.high}). The
engine-derived probability is overconfident at both extremes
(Figure~\ref{fig:practical}): where it gives the mover 80 to 90\%,
the mean result is \V{practical-chances.slice.win.80-90.happened}. The error
is largest for lower-rated players: from positions rated at 80\% or more,
players rated under 1100 score
\V{practical-chances.grid.rating-win.-1100.80+.happened} and players rated
2000 and over score
\V{practical-chances.grid.rating-win.2000+.80+.happened}. It also ignores
the clock: with under thirty seconds remaining the mean result is
\V{practical-chances.slice.clock.-30.happened}, where the engine-derived
probability averages
\V{practical-chances.slice.clock.-30.engine-curve.said}. The best network
is not significantly different from the GBDT baseline in this task, and
pretraining makes the transformers worse (Section~\ref{sec:start}). The best
model is \V{practical-chances.best.name}, the GBDT baseline trained on four
times the standard training set.

\begin{figure}[htbp]\centering
\block{practical-chances-reliability-engine-curve}
\caption{Practical chances: mean game result against the engine-derived win
probability, in equal-sized bins of the test set ordered by that
probability, with the GBDT baseline for comparison.}
\label{fig:practical}
\end{figure}

\subsection{Where to look}

A model ranks the squares of each test position by predicted probability,
and the metric is the average precision of the ranking. The GBDT baseline
scores \V{where-to-look.reference.score}, the best network
(\V{where-to-look.network.name}) \V{where-to-look.network.score} and the
constant baseline \V{where-to-look.floor.score}. The top-ranked square is
part of the solution in \V{where-to-look.reference.top1.percent}\% of
positions for the GBDT baseline, \V{where-to-look.network.top1.percent}\%
for the network and \V{where-to-look.floor.top1.percent}\% for the constant
baseline. Table~\ref{tab:look} gives the average precision by puzzle rating.

\begin{table}[htbp]\centering\small
\block{where-to-look-bands-table}
\caption{Where to look: average precision of each model by the puzzle's
rating.}
\label{tab:look}
\end{table}

Identifying one relevant square is the easy part of the task: GBDT on the
context features alone, which include whether the opponent's last move
started or ended on the square, ranks a solution square first in
\V{where-to-look.run.gbm-context.top1.percent}\% of positions. Ranking all
solution squares ahead of the others is harder, and it is there that the
network leads. The task also becomes harder with the puzzle's rating: the
network's average precision falls from \V{where-to-look.band.1.network} for
puzzles rated \V{where-to-look.band.1.name} to
\V{where-to-look.band.5.network} for puzzles rated
\V{where-to-look.band.5.name}. The predicted probabilities are calibrated to
within \V{where-to-look.best.calibration.widest-gap} percentage points in
each of \V{where-to-look.best.calibration.groups} bins. Every position in
the dataset contains a tactic, so the model predicts where a known tactic is
located and not whether one exists.

\subsection{Puzzle themes}

A model ranks the \V{puzzle-themes.task.candidates.count} theme tags for
each test puzzle, and the metric is the average precision of the ranking
over the puzzles that carry at least one tag. The GBDT baseline scores
\V{puzzle-themes.reference.score}, the best network
(\V{puzzle-themes.network.name}) \V{puzzle-themes.network.score} and the
constant baseline \V{puzzle-themes.floor.score}. The top-ranked tag is
correct in \V{puzzle-themes.reference.top1.percent}\% of puzzles for the GBDT
baseline, \V{puzzle-themes.network.top1.percent}\% for the network and
\V{puzzle-themes.floor.top1.percent}\% for the constant baseline.

An average over puzzles is dominated by the common tags, so
Table~\ref{tab:themes} evaluates each tag separately as a binary
classification. The most common tag is
\V{puzzle-themes.candidates.common.name}
(\V{puzzle-themes.candidates.common.share}\% of test puzzles) and the rarest
is \V{puzzle-themes.candidates.rare.name}
(\V{puzzle-themes.candidates.rare.share}\%). For
\V{puzzle-themes.candidates.run}, the tag with the highest ROC-AUC is
\V{puzzle-themes.candidates.roc.most.name}
(\V{puzzle-themes.candidates.roc.most}) and the tag with the lowest is
\V{puzzle-themes.candidates.roc.least.name}
(\V{puzzle-themes.candidates.roc.least}).

\begin{table}[htbp]\centering\small
\block{puzzle-themes-candidates}
\caption{Puzzle themes: each tag evaluated separately on the test set, most
common first. \emph{Test examples} and \emph{Share} give the number and
percentage of test puzzles that carry the tag. PR-AUC and ROC-AUC measure how
well a model's probability for the tag separates those puzzles from the
rest.}
\label{tab:themes}
\end{table}

By ROC-AUC the network separates every tag almost perfectly. PR-AUC is the
stricter measure, because it penalises false positives on rare tags. By
PR-AUC the \V{puzzle-themes.candidates.pr.weakest.count} weakest tags are
\V{puzzle-themes.candidates.pr.weakest.names}, the lowest at
\V{puzzle-themes.candidates.pr.least} and all at or below
\V{puzzle-themes.candidates.pr.weakest.highest}. Rarity does not explain the
weak tags: the tag with the highest PR-AUC,
\V{puzzle-themes.candidates.pr.most.name}, appears in
\V{puzzle-themes.candidates.pr.most.share}\% of test puzzles, and the tag
with the lowest in \V{puzzle-themes.candidates.pr.least.share}\%. We
hypothesise that the well-predicted tags correspond to a single clear
pattern on the board, whereas the poorly predicted tags depend on what the
solution does not show, such as a piece with no safe square or a side with
no good move. The network's advantage over GBDT is smallest for
\V{puzzle-themes.candidates.pr.lead.least.name}
(\V{puzzle-themes.candidates.pr.lead.least.reference} for GBDT and
\V{puzzle-themes.candidates.pr.lead.least.best} for the network) and largest
for \V{puzzle-themes.candidates.pr.lead.most.name}
(\V{puzzle-themes.candidates.pr.lead.most.reference} and
\V{puzzle-themes.candidates.pr.lead.most.best}). The tags are produced by
Lichess's own heuristics and user votes, so a model that predicts them
reproduces that labelling, including its errors.

\subsection{Human move}
\label{sec:human-move}

A test position has \V{human-move.candidates.mean} legal moves on average.
Maia-2 as released has a top-1 accuracy of \V{human-move.run.maia2.top1} and
a top-3 accuracy of \V{human-move.run.maia2.top3}. The best network trained
here, \V{human-move.network.name}, reaches \V{human-move.network.top1} and
\V{human-move.network.top3}. The GBDT baseline reaches
\V{human-move.reference.top1} top-1 accuracy and the uniform baseline
\V{human-move.floor.top1}. Table~\ref{tab:move} gives top-1 accuracy by the
mover's rating.

\begin{table}[htbp]\centering\small
\block{human-move-bands-table}
\caption{Human move: top-1 accuracy of each model by the mover's rating.}
\label{tab:move}
\end{table}

That Maia-2 leads is expected, since it was trained on years of games and
our networks on the games of one month. The evaluation adds three
observations.

\paragraph{Maia-2 remains calibrated.}
The test games were played years after Maia-2's training data. The only
parameter we fit is a softmax temperature on the validation set, and its
value is \V{human-move.run.maia2.temperature}, where 1 leaves the released
probabilities unchanged. With it, none of the
\V{human-move.best.calibration.groups} probability bins deviates from the
observed frequency by more than \V{human-move.best.calibration.widest-gap}
percentage points.

\paragraph{Rating conditioning matters at the extremes.}
When Maia-2 is given the same fixed rating for both players in every
position (\V{human-move.run.maia2-one-rating.name}), its top-1 accuracy
falls to \V{human-move.run.maia2-one-rating.top1}, a paired gain of
\V{human-move.pair.maia2-one-rating.maia2.gain} (95\% CI
\V{human-move.pair.maia2-one-rating.maia2.gain.low} to
\V{human-move.pair.maia2-one-rating.maia2.gain.high};
Table~\ref{tab:move-studies}). In the band that contains the fixed rating,
\V{human-move.band.3.name}, accuracy is almost unchanged
(\V{human-move.run.maia2.band.3} and
\V{human-move.run.maia2-one-rating.band.3}). The loss is largest at the two
ends: from \V{human-move.run.maia2.band.1} to
\V{human-move.run.maia2-one-rating.band.1} for players rated
\V{human-move.band.1.name}, and from \V{human-move.run.maia2.band.5} to
\V{human-move.run.maia2-one-rating.band.5} for players rated
\V{human-move.band.5.name}.

\begin{table}[htbp]\centering\small
\block{human-move-studies-bands}
\caption{Human move: top-1 accuracy of each ablation model, listed under the
model it modifies, by the mover's rating.}
\label{tab:move-studies}
\end{table}

\paragraph{Our networks are limited by data and sensitive to architecture.}
For the lowest-rated players \V{human-move.network.name} is close to Maia-2
(\V{human-move.band.1.network} compared with
\V{human-move.run.maia2.band.1}); for the highest-rated it is far behind
(\V{human-move.band.5.network} compared with
\V{human-move.run.maia2.band.5}). Model size has little effect: transformers
with \V{human-move.run.transformer-small.parameters.millions},
\V{human-move.run.transformer-medium.parameters.millions} and
\V{human-move.run.transformer-large.parameters.millions} million parameters
reach \V{human-move.run.transformer-small.top1},
\V{human-move.run.transformer-medium.top1} and
\V{human-move.run.transformer-large.top1}, and the paired gain of the largest
over the medium one is
\V{human-move.pair.transformer-large.transformer-medium.gain}. Data has a
larger effect: the paired gain of the medium transformer is
\V{human-move.pair.transformer-medium-half.transformer-medium.gain} on half
of the training games and
\V{human-move.pair.transformer-medium-quarter.transformer-medium.gain} on a
quarter. Among the architecture modifications, the attack-aware attention
bias gains
\V{human-move.pair.transformer-medium-attacks.transformer-medium.gain} (95\%
CI \V{human-move.pair.transformer-medium-attacks.transformer-medium.gain.low}
to \V{human-move.pair.transformer-medium-attacks.transformer-medium.gain.high})
and is accepted. The ResNet and the hybrid are significantly worse (gains of
\V{human-move.pair.conv-medium.transformer-medium.gain} and
\V{human-move.pair.hybrid-medium.transformer-medium.gain}). Removing the
relative position bias has a gain of
\V{human-move.pair.transformer-medium-no-geometry.transformer-medium.gain},
which leaves the transformer level with the GBDT baseline (gain over that
baseline \V{human-move.run.transformer-medium-no-geometry.gain}, 95\% CI
\V{human-move.run.transformer-medium-no-geometry.gain.low} to
\V{human-move.run.transformer-medium-no-geometry.gain.high}). Adding the
hand-crafted features is
\V{human-move.pair.transformer-medium-hand.transformer-medium.verdict}. A
plausible explanation for the architecture results is that a move connects
two squares that may be far apart: attention relates any two squares in one
layer, whereas a convolution relates only neighbours, and to do so attention
needs to know how the two squares are positioned relative to each other and
whether the piece on one reaches the other.

\subsection{Game rating}
\label{sec:game-rating}

An example is one player in one game, and the target is that player's rating
at the time. The GBDT baseline has an MAE of
\V{game-rating.reference.score} rating points, the best network
(\V{game-rating.network.name}) \V{game-rating.network.score} and the
constant baseline \V{game-rating.floor.score}. GBDT on the time control and
the length of the game alone has an MAE of
\V{game-rating.run.gbm-context.score}. The best model's predictions have a
Spearman rank correlation of \V{game-rating.best.spearman} with the true
ratings. Table~\ref{tab:rating} gives the error by the player's rating. The
bands are defined on the target itself, so a model that regresses towards
the mean has larger errors at both ends.

\begin{table}[htbp]\centering\small
\block{game-rating-bands-table}
\caption{Game rating: MAE of each model in rating points, by the player's
rating.}
\label{tab:rating}
\end{table}

In this task \V{game-rating.ladder.ahead.of} are significantly better than
the GBDT baseline, and the best is the smallest, with
\V{game-rating.network.parameters.millions} million parameters compared
with \V{game-rating.run.transformer-medium.parameters.millions} million for
the medium transformer. Adding the hand-crafted features makes the medium
transformer worse (paired gain
\V{game-rating.pair.transformer-medium-hand.transformer-medium.gain} rating
points, 95\% CI
\V{game-rating.pair.transformer-medium-hand.transformer-medium.gain.low} to
\V{game-rating.pair.transformer-medium-hand.transformer-medium.gain.high}),
the opposite of most other tasks (Table~\ref{tab:changes}).

A prediction from a single game is imprecise: the best model is within 200
points of the player's rating for
\V{game-rating.best.close.within.200}\% of test examples. The prediction
also depends on the opponent. Among players rated 1400 to 1700, the best
network predicts \V{game-rating.grid.rating-gap.1400-1700.--200.said} on average against opponents rated 200 or
more points higher and \V{game-rating.grid.rating-gap.1400-1700.200+.said} against opponents rated 200 or more
points lower, although the players' mean ratings in the two groups are
\V{game-rating.grid.rating-gap.1400-1700.--200.happened} and \V{game-rating.grid.rating-gap.1400-1700.200+.happened}. The GBDT baseline, which
receives features of the player's own moves only, shows the same dependence
(\V{game-rating.grid.rating-gap.1400-1700.--200.gbm-hand.said} and \V{game-rating.grid.rating-gap.1400-1700.200+.gbm-hand.said}), and GBDT on
the context features does not (\V{game-rating.grid.rating-gap.1400-1700.--200.gbm-context.said} and
\V{game-rating.grid.rating-gap.1400-1700.200+.gbm-context.said}). A plausible explanation is that a weaker
opponent creates more opportunities, and each missed opportunity is counted
as a loss of accuracy.

\section{Practical Use}
\label{sec:use}

The models were built for a coaching application. The results suggest the
following for deployment.

\begin{itemize}
\item \textbf{Start with GBDT.} Where the models see every input, no
network gains more than \V{lab.network.gain.whole.most} points of normalised
score over GBDT, and the GBDT models train in minutes and run without a GPU.
A network (\V{puzzle-difficulty.best.weights.megabytes} MB of weights for
the best puzzle-difficulty model) is worth deploying only where its
advantage matters for the application.
\item \textbf{Use a lookup table where context dominates.} In think time,
blunder hazard and practical chances most of the normalised score comes from
ratings, clocks and, where available, the engine evaluation. A table over
these inputs is a reasonable first version with no model.
\item \textbf{Calibrated probabilities can be shown.} Blunder hazard and
tactic found produce probabilities that match observed frequencies to within
\V{blunder-hazard.best.calibration.widest-gap} and
\V{tactic-found.best.calibration.widest-gap} percentage points in every
bin. They describe players of a given rating in general, and should be
presented that way and not as a statement about the individual.
\item \textbf{Puzzle difficulty supports ranking, not rating.} The model can
order the tactics of a game from routine to hard. Its error on a single
tactic is too large to display a rating.
\item \textbf{Think time supports patterns, not single moves.} The model
gives the typical time for a position, rating and clock. A review can count
how often a player moved much faster than typical and in what kind of
position, but should not label one move as too fast.
\item \textbf{Practical chances complements the engine evaluation.} It gives
the expected result between two players of given ratings and clocks, next to
the result under best play. Positions where the two differ strongly are
candidates for a coach's attention.
\item \textbf{Rate a tactic from its solution.} The solution is worth
\V{puzzle-difficulty-unsolved.shown.best.worth} rating points of RMSE
between the best models of the two puzzle-difficulty tasks. An application
that has the engine's line should therefore estimate difficulty from the
line.
\item \textbf{Where to look and puzzle themes add hints and names.} For a
tactic that an engine has found and the player missed, where to look
indicates the relevant squares before the move is revealed, and puzzle
themes names the motif for the tags that Table~\ref{tab:themes} shows are
predicted well.
\item \textbf{Maia-2 can be used as released.} Its probabilities remained
calibrated on these games. The ratings it is given matter most for the
lowest- and highest-rated players.
\item \textbf{Game rating is a wide estimate.} It can describe what one game
looked like for games between players of similar rating. It cannot yet say
whether a player performed above or below their rating, because against a
much stronger or weaker opponent the prediction mostly reflects the
opponent.
\end{itemize}

All results hold for Lichess players and Lichess ratings. Ratings from
another site must first be mapped to that scale, and any model that is
deployed should be validated on the target player's own games.

\section{Limitations}
\label{sec:limits}

\begin{itemize}
\item \textbf{One site, one month, one time control.} The game positions
come from rated \V{blunder-hazard.recipe.speeds} games played on Lichess in
\V{blunder-hazard.data.month}. We do not know how the results transfer to
other time controls or sites.
\item \textbf{Selection of analysed games.} A position has an engine
evaluation only if someone requested server analysis of its game, and
players choose which games to analyse.
\item \textbf{Shared test games.} Blunder hazard, think time, human move and
practical chances are evaluated on the same
\V{blunder-hazard.data.games.test} test games, so their results are not
independent, and the confidence intervals are as wide as that number of
games allows. Game rating uses games of the same month under the same split.
\item \textbf{Selection of puzzles.} The puzzle tasks contain only tactics
that Lichess turned into puzzles, which may differ from tactics in general.
\item \textbf{Small models and datasets.} The largest network has
\V{lab.networks.largest.parameters.millions} million parameters, and the
largest task-specific training set has
\V{puzzle-difficulty.run.gbm-hand-quadruple.train-rows} examples
(\V{puzzle-difficulty.data.rows.train} in the standard training set). The
comparison between
networks and GBDT with hand-crafted features is a statement about models of
this size on this amount of data.
\item \textbf{A single seed.} Every network was trained once. The
confidence intervals cover sampling of the test games and not the variation
between training runs, so gains of a few points of normalised score need
confirmation with additional seeds.
\item \textbf{A single pretrained encoder.} Pretraining was evaluated with
one encoder size, one schedule, one sample of evaluations and one
pretraining target.
\item \textbf{The normalised score ignores the uncertainty of its
denominator.} It is intended for comparison across tasks. Conclusions about
one task should rest on the paired gains in that task's own metric.
\end{itemize}

\section{Conclusion and Future Work}
\label{sec:next}

Across \V{lab.side.count} tasks about human play, GBDT with hand-crafted
features is a strong and inexpensive baseline, and the best neural network
improves on it significantly in \V{lab.network.ahead.of}, usually by a
modest margin. The most reliable way to improve a network was to pretrain
its encoder on engine evaluations. Architecture changes had smaller and
less consistent effects, and more data helped every model class.

Several experiments follow directly from these results. The first is to
train with more data and more optimisation steps together, which neither
data-size ablation covers. The second is to combine pretraining with the
modifications that were accepted when trained from scratch, in particular
the attack-aware attention bias, and to evaluate pretraining with a longer
schedule, a larger encoder, and a pretraining target based on human moves
instead of engine evaluations. The third is to repeat the main comparisons
with several random seeds, so that the smaller gains can be confirmed or
rejected. Finally, where to look, puzzle themes, human move and game rating
have so far received only the fixed sequence of models and a subset of the
modifications. Pretraining has not been evaluated on them, and where to look
is a natural candidate because, like the unsolved puzzle task, it withholds
the solution.

\bibliographystyle{unsrt}
\bibliography{../references}

\appendix

\section{Results for Every Model}
\label{sec:runs}

Each table lists the models of one task ranked by the primary metric on the
test set, with the 95\% CI of the score and the paired gain over the GBDT
baseline with its 95\% CI. Gains in bold are statistically significant. A
table holds the fixed sequence of models, the formula baselines and every
modification that was accepted; the task papers also list the modifications
that were not. In a model's identifier, \texttt{-hand} adds the hand-crafted
features, \texttt{-double} and \texttt{-quadruple} train on twice and four
times the standard training set, and the other suffixes are defined in
Section~\ref{sec:modifications}.

\begin{table}[htbp]\centering\small
\block{puzzle-difficulty-leaderboard}
\caption{Puzzle difficulty.}
\end{table}

\begin{table}[htbp]\centering\small
\block{puzzle-difficulty-unsolved-leaderboard}
\caption{Puzzle difficulty unsolved.}
\end{table}

\begin{table}[htbp]\centering\small
\block{blunder-hazard-leaderboard}
\caption{Blunder hazard.}
\end{table}

\begin{table}[htbp]\centering\small
\block{think-time-leaderboard}
\caption{Think time.}
\end{table}

\begin{table}[htbp]\centering\small
\block{tactic-found-leaderboard}
\caption{Tactic found.}
\end{table}

\begin{table}[htbp]\centering\small
\block{practical-chances-leaderboard}
\caption{Practical chances.}
\end{table}

\begin{table}[htbp]\centering\small
\block{where-to-look-leaderboard}
\caption{Where to look.}
\end{table}

\begin{table}[htbp]\centering\small
\block{puzzle-themes-leaderboard}
\caption{Puzzle themes.}
\end{table}

\begin{table}[htbp]\centering\small
\block{human-move-leaderboard}
\caption{Human move.}
\end{table}

\begin{table}[htbp]\centering\small
\block{game-rating-leaderboard}
\caption{Game rating.}
\end{table}

\begin{table}[htbp]\centering\small
\block{position-evaluation-leaderboard}
\caption{Position evaluation, the pretraining task.}
\end{table}

\end{document}
